\documentclass[12pt]{article}
\usepackage{amsmath,amsthm,amssymb,mathtools}
\usepackage{hyperref}

\newtheorem{theorem}{Theorem}[section]
\newtheorem{lemma}[theorem]{Lemma}
\newtheorem{proposition}[theorem]{Proposition}
\newtheorem{definition}[theorem]{Definition}
\newtheorem{remark}[theorem]{Remark}

\title{Mollifier Construction for Symmetric Square $L$-Functions\\
\large \texttt{M-1, M-2 Deliverables --- sym2-effective-lbound}}
\author{}
\date{Status: \texttt{[OBL]} --- proof obligations open}

\begin{document}
\maketitle

\begin{abstract}
This document lays out the mollifier construction for $L(s,\mathrm{sym}^2 f)$,
following the Goldfeld--Hoffstein--Lieman approach. We define an explicit mollifier
$M(s)$ of length $X = N^\theta$, state the required mean value theorem (M-2),
and connect nonvanishing at $s=1$ to the effective lower bound goal (E-1).
The proof of M-2 requires the GL$_3$ Rankin--Selberg mean value theorem: \texttt{[OBL]}.
\end{abstract}

\tableofcontents

\section{The Symmetric Square as a GL$_3$ Object}

\begin{theorem}[Gelbart--Jacquet 1978, \texttt{[BASE]}]\label{thm:GJ}
For a cuspidal $\pi$ of $GL_2(\mathbb{A})$, there exists a cuspidal or isobaric
$\Pi = \mathrm{sym}^2(\pi)$ of $GL_3(\mathbb{A})$ with $L(s,\Pi) = L(s,\mathrm{sym}^2 f)$.
\end{theorem}

The Dirichlet series is $L(s,\mathrm{sym}^2 f) = \sum_{n=1}^\infty a_{\mathrm{sym}^2}(n)/n^s$,
where $a_{\mathrm{sym}^2}(p) = \tilde\alpha_p^2 + 1 + \tilde\beta_p^2 = a_f(p)^2/p^{k-1} - 1$
at unramified $p$.

\section{The Mollifier}

\begin{definition}\label{def:mollifier}
For length parameter $X \geq 1$:
\[
  M(s) = \sum_{n \leq X} \frac{\mu(n)\,a_{\mathrm{sym}^2}(n)}{n^s}.
\]
\end{definition}

The heuristic: $M(s)\cdot L(s,\mathrm{sym}^2 f) \approx 1$ near $s=1$
since $M(s) \approx L(s,\mathrm{sym}^2 f)^{-1}$ via M\"obius inversion.

\section{Mean Value Theorem \texttt{[OBL]}}

\begin{theorem}[M-2, \texttt{[OBL]}]\label{thm:M2}
For $T \geq 1$ and $X = N^\theta$ ($\theta > 0$ to be optimized):
\[
  \int_T^{2T} \bigl|M(\tfrac12+it)\,L(\tfrac12+it,\mathrm{sym}^2 f)\bigr|^2\,dt \gg T.
\]
The implied constant is effective, depending on $\theta$, $k$, and $N$.
\end{theorem}

\begin{remark}[Proof strategy]
Requires: (a) GL$_3$ Rankin--Selberg mean value theorem via Theorem~\ref{thm:GJ};
(b) mollifier mean value via GL$_3$ large sieve;
(c) parameter optimization.
The main gap is the GL$_3$ large sieve with an explicit constant: \texttt{[OBL]}.
\end{remark}

\section{Connection to the Effective Bound (E-1)}

The effective lower bound $L(1,\mathrm{sym}^2 f) \geq c_{\mathrm{eff}}/\log N$ follows from:
\begin{enumerate}
  \item[Step 1] Zero-free region (M-3): $L(s,\mathrm{sym}^2 f)\neq 0$ for
    $\mathrm{Re}(s) > 1 - c_0/\log N$ with explicit $c_0$.
  \item[Step 2] Contour integral: zero-free region $\Rightarrow$
    $L(1,\cdot) \geq c_1/\log N$ with $c_1$ explicitly computable from $c_0$.
  \item[Step 3] Mollifier (M-2) certifies Case~2 (Siegel zero) cannot arise.
\end{enumerate}
The explicit value of $c_1$ is the core research deliverable (E-1).

\end{document}
